% Part:sets-functions-relations
% Chapter: size-of-sets
% Section: enumerations

\documentclass[../../../include/open-logic-section]{subfiles}

\begin{document}

\olfileid[fa-IR]{sfr}{siz}{enm}

\olsection{برشماری‌ها و مجموعه‌های \usetoken{S}{enumerable}}

\begin{editorial}
  این بخش برشماری‌های مجموعه‌ها را بررسی می‌کند و آن‌ها را توابع
  پوشا از $\PosInt$ تعریف می‌کند. مطالب را برای خوانندگانی که پیشینهٔ
  ریاضی اندکی دارند، آهسته و گام‌به‌گام پیش می‌برد. روایتی جایگزین و
  موجزتر در \olref[enm-alt]{sec} آمده است که برشماری‌ها را به‌گونه‌ای
  متفاوت تعریف می‌کند: به‌منزلهٔ تناظرهای دوسویی با $\Nat$ (یا یک قطعهٔ آغازین).
\end{editorial}

\begin{explain}
پیش‌تر نمونه‌هایی از مجموعه‌ها را با فهرست‌کردن !!{element}sی آن‌ها
ارائه کرده‌ایم. بیایید به‌طور کلی‌تر بررسی کنیم که چگونه و چه هنگام
می‌توانیم !!{element}sی یک مجموعه را فهرست کنیم، حتی اگر آن مجموعه
نامتناهی باشد.
\end{explain}

\begin{defn}[برشماری، به‌طور غیرصوری]
به‌طور غیرصوری، یک \emph{برشماری} مجموعهٔ~$A$ فهرستی (احتمالاً
نامتناهی) از !!{element}sی~$A$ است، به‌گونه‌ای که هر !!{element}
از $A$ در جایگاهی متناهی از فهرست ظاهر شود. اگر $A$ برشماری‌ای داشته
باشد، می‌گوییم $A$ \emph{!!{enumerable}} است.
\end{defn}

\begin{explain}
چند نکته دربارهٔ برشماری‌ها:
\begin{enumerate}
\item تنها فهرست‌هایی را برشماری به شمار می‌آوریم که آغازی داشته باشند و
  در آن‌ها هر !!{element} جز نخستین عضو، دقیقاً یک
  !!{element} داشته باشد که بی‌درنگ پیش از آن بیاید. به بیان دیگر، میان
  نخستین !!{element} آن فهرست و هر !!{element} دیگر فقط شمار متناهی‌ای عضو
  قرار دارد. به‌ویژه، این یعنی هر !!{element} در یک برشماری جایگاهی متناهی
  دارد: نخستین !!{element} در جایگاه~$1$، دومین در جایگاه~$2$ و الی آخر است.
\item می‌توانیم برشماری‌های متفاوتی از همان مجموعهٔ~$A$ داشته باشیم که
  در ترتیب ظاهرشدن !!{element}s فرق کنند: $4$، $1$،
  $25$، $16$،~$9$ پنج عدد مربعی نخست را (به‌منزلهٔ یک مجموعه) درست
  به همان خوبیِ $1$، $4$، $9$، $16$،~$25$ برمی‌شمارد.
\item برشماری‌های تکراردار همچنان برشماری‌اند: $1$، $1$، $2$،
  $2$، $3$، $3$،~\dots{} همان مجموعه‌ای را برمی‌شمارد که $1$، $2$،
  $3$،~\dots{} برمی‌شمارد.
\item وقتی برشماری‌ای را مشخص می‌کنیم، ترتیب و تکرار \emph{اهمیت
  دارند}: می‌توانیم برشماریِ اعداد صحیح مثبت را با $1$، $2$، $3$، $1$،
  \dots{} آغاز کنیم، اما وقتی آن‌ها را به شیوهٔ متعارف، یعنی $1$، $2$،
  $3$، $4$،~\dots، برمی‌شماریم، الگو آسان‌تر دیده می‌شود.
\item برشماری‌ها باید آغاز داشته باشند: \dots، $3$، $2$، $1$ برشماریِ
  اعداد صحیح مثبت نیست، زیرا نخستین !!{element} ندارد. برای اینکه ببینید
  این نکته چگونه از تعریف غیرصوری نتیجه می‌شود، از خود بپرسید: ``عدد 76
  در کدام جایگاه فهرست ظاهر می‌شود؟''
\item عبارت زیر برشماریِ اعداد صحیح مثبت نیست:
  $1$، $3$، $5$، \dots، $2$، $4$، $6$، \dots\@ مشکل این است که
  اعداد زوج به‌جای جایگاه‌های متناهی، در جایگاه‌های $\infty + 1$،
  $\infty + 2$، $\infty +
  3$ قرار می‌گیرند.
\item مجموعهٔ تهی شمارا است: فهرست تهی آن را برمی‌شمارد!{}
\end{enumerate}
\end{explain}

\begin{prop}
  اگر $A$ برشماری‌ای داشته باشد، برشماری‌ای بی‌تکرار نیز دارد.
\end{prop}

\begin{proof}
  فرض کنید $A$ دارای برشماریِ $x_1$، $x_2$، \dots{} باشد که در آن هر
  $x_i$ یک !!{element} از~$A$ است. می‌توانیم با حذف !!{element}sی تکراری،
  تکرارها را از یک برشماری برداریم. برای نمونه، می‌توانیم برشماری را به
  برشماری تازه‌ای تبدیل کنیم که در آن $x_i$ را، اگر !!a{element} از~$A$
  باشد که در میان $x_1$، \dots، $x_{i-1}$ نیست، وارد فهرست کنیم؛ و اگر
  $x_i$ پیش‌تر در میان $x_1$، \dots،~$x_{i-1}$ ظاهر شده است، آن را از
  فهرست حذف کنیم.
\end{proof}

استدلال اخیر نشان می‌دهد که برای دستیابی به درکی دقیق از
برشماری‌ها و مجموعه‌های !!{enumerable} و اثبات مطالبی دربارهٔ آن‌ها،
به تعریفی دقیق‌تر نیاز داریم. تعریف زیر چنین تعریفی را فراهم می‌کند.

\begin{defn}[برشماری، به‌طور صوری]
منظور از یک \emph{برشماری} برای مجموعهٔ $A \neq \emptyset$، هر تابع
!!{surjective} $f \colon \PosInt \to A$ است.
\end{defn}

\begin{explain}
بیایید خود را متقاعد کنیم که تعریف صوری و تعریف غیرصوری با استفاده از
فهرستی احتمالاً نامتناهی، هم‌ارزند. نخست، هر تابع !!{surjective} از
$\PosInt$ به مجموعهٔ~$A$، مجموعهٔ~$A$ را برمی‌شمارد. چنین تابعی
برشماری‌ای به معنای غیرصوری بالا تعیین می‌کند: فهرست $f(1)$، $f(2)$،
$f(3)$، \dots. چون $f$ !!{surjective} است، تضمین می‌شود که هر
!!{element} از~$A$ مقدارِ~$f(n)$ برای یک~$n \in \PosInt$ باشد. ازاین‌رو،
هر !!{element} از $A$ در جایگاهی متناهی از فهرست ظاهر می‌شود. چون ممکن
است تابع !!{injective} نباشد، فهرست می‌تواند تکراردار باشد، اما این
پذیرفتنی است (چنان‌که در بالا گفته شد).

از سوی دیگر، با داشتن فهرستی که همهٔ !!{element}sی~$A$ را
برمی‌شمارد، می‌توانیم تابعِ !!a{surjective} $f\colon \PosInt \to A$ را
چنان تعریف کنیم که $f(n)$ برابر $n$اُمین !!{element} در فهرست باشد، یا اگر
در جایگاه $n$ هیچ !!{element}ی نیست، برابر آخرین !!{element} در فهرست باشد.
تنها حالتی که این کار تابعی !!a{surjective} تولید نمی‌کند زمانی است که
$A$ تهی و در نتیجه فهرست تهی باشد. پس هر فهرست ناتهی تابعی
!!a{surjective} به‌صورت $f\colon \PosInt \to A$ تعیین می‌کند.
\end{explain}

\begin{defn}
  \ollabel{defn:enumerable}
  مجموعهٔ~$A$ !!{enumerable} است اگر و تنها اگر تهی باشد یا برشماری‌ای داشته باشد.
\end{defn}

\begin{ex}
تابعی که اعداد صحیح مثبت ($\PosInt$) را برمی‌شمارد، صرفاً تابع همانیِ
$f(n) = n$ است. تابعی که اعداد طبیعی $\Nat$ را برمی‌شمارد، تابع
$g(n) = n - 1$ است.
\end{ex}

\begin{ex}
توابع $f\colon \PosInt \to \PosInt$ و $g \colon \PosInt \to
\PosInt$ که به‌صورت زیر داده می‌شوند
\begin{align*}
f(n) & = 2n \text{ و}\\
g(n) & = 2n - 1
\end{align*}
به‌ترتیب اعداد صحیح مثبت زوج و اعداد صحیح مثبت فرد را برمی‌شمارند.
بااین‌حال، هیچ‌یک از این دو تابع برشماریِ $\PosInt$ نیست، زیرا هیچ‌یک
!!{surjective} نیست.
\end{ex}

\begin{prob}
برشماری‌ای از مربع‌های مثبت $1$، $4$، $9$، $16$، \dots تعریف کنید.
\end{prob}

\begin{ex}
تابع $f(n) = (-1)^{n} \lceil \frac{(n-1)}{2}\rceil$ (که در آن
$\lceil x \rceil$ تابع \emph{سقف} را نشان می‌دهد که $x$ را رو به بالا
تا نزدیک‌ترین عدد صحیح گرد می‌کند) مجموعهٔ اعداد صحیح~$\Int$ را
برمی‌شمارد. توجه کنید که $f$ چگونه مقادیر $\Int$ را با ``جهیدن'' رفت‌وبرگشتی
میان اعداد صحیح مثبت و منفی تولید می‌کند:
\[
\begin{array}{c c c c c c c c}
f(1) & f(2) & f(3) & f(4) & f(5) & f(6) & f(7) & \dots \\ \\
- \lceil \tfrac{0}{2} \rceil & \lceil \tfrac{1}{2}\rceil & - \lceil \tfrac{2}{2} \rceil & \lceil \tfrac{3}{2} \rceil & - \lceil \tfrac{4}{2} \rceil  & \lceil \tfrac{5}{2}
\rceil & - \lceil \tfrac{6}{2} \rceil & \dots \\ \\
0 & 1 & -1 & 2 & -2 & 3 & \dots
\end{array}
\]
همچنین می‌توانید $f$ را به‌صورت موردیِ زیر در نظر بگیرید:
\[
f(n) = \begin{cases}
  0 & \text{اگر $n = 1$}\\
  n/2 & \text{اگر $n$ زوج باشد}\\
  -(n-1)/2 & \text{اگر $n$ فرد و $>1$ باشد}
  \end{cases}
\]
\end{ex}

\begin{prob}
  نشان دهید اگر $A$ و $B$ !!{enumerable} باشند، $A \cup B$ نیز چنین
  است. برای این کار، فرض کنید توابع !!{surjective} $f\colon \PosInt \to
  A$ و $g\colon \PosInt \to B$ وجود دارند، و تابعی
  !!a{surjective}~$h\colon \PosInt \to A \cup B$ تعریف کنید و اثبات کنید که
  !!{surjective} است. حالت‌هایی را که $A$ یا~$B = \emptyset$ است نیز بررسی کنید.
\end{prob}

\begin{prob}
  نشان دهید اگر $B \subseteq A$ و $A$ !!{enumerable} باشد، $B$ نیز چنین
  است. برای این کار، فرض کنید تابعی !!a{surjective} به‌صورت $f\colon \PosInt \to
  A$ وجود دارد. تابعی !!a{surjective} به‌صورت $g\colon \PosInt \to B$
  تعریف کنید و اثبات کنید که !!{surjective} است. اگر $B = \emptyset$ باشد چه رخ می‌دهد؟
\end{prob}

\begin{prob}
  با استقرا بر $n$ نشان دهید که اگر $A_1$، $A_2$، \dots، $A_n$ همگی
  !!{enumerable} باشند، $A_1 \cup \dots \cup A_n$ نیز چنین است. می‌توانید
  این واقعیت را فرض کنید که اگر دو مجموعهٔ $A$ و~$B$ !!{enumerable}
  باشند، $A \cup
  B$ نیز چنین است.
\end{prob}

هرچند هنگام فهرست‌کردن !!{element}sی یک مجموعه شاید طبیعی‌تر باشد که
شمردن را از نخستین !!{element}، یعنی جایگاه $1$، آغاز کنیم، ریاضی‌دانان
دوست دارند برای شمردن چیزها از اعداد طبیعی~$\Nat$ استفاده کنند. آنان از
!!{element}sی $0$اُم، $1$اُم، $2$اُم و الی آخرِ یک فهرست سخن
می‌گویند. بر همین اساس، می‌توانیم برشماری را تابعی !!a{surjective} از
$\Nat$ به~$A$ تعریف کنیم. البته این دو تعریف هم‌ارزند.

\begin{prop}\ollabel{prop:enum-shift}
  !!a{surjection} $f\colon \PosInt \to A$ وجود دارد اگر و تنها اگر
  !!a{surjection} $g\colon \Nat \to A$ وجود داشته باشد.
\end{prop}

\begin{proof}
  با داشتن !!a{surjection} $f\colon \PosInt \to A$، می‌توانیم تعریف کنیم
  $g(n) =
  f(n+1)$ برای همهٔ $n \in \Nat$. به‌آسانی می‌توان دید که $g\colon \Nat
  \to A$ !!{surjective} است. برعکس، با داشتن !!a{surjection} $g\colon
  \Nat \to A$، تعریف کنید $f(n) = g(n-1)$.
\end{proof}

این امر نتیجهٔ زیر را به ما می‌دهد:

\begin{cor}\ollabel{cor:enum-nat}
مجموعهٔ $A$ !!{enumerable} است اگر و تنها اگر تهی باشد یا تابعی
!!a{surjective} به‌صورت $f\colon \Nat \to A$ وجود داشته باشد.
\end{cor}

در بالا گفتیم که می‌توان فهرستی از !!{element}sی مجموعهٔ~$A$ را به
فهرستی بی‌تکرار تبدیل کرد. این مطلب دربارهٔ برشماری‌ها نیز درست است،
اما صورت‌بندی و اثبات دقیق آن کمی دشوارتر است. هر تابع $f\colon \PosInt \to A$
باید برای همهٔ $n \in
\PosInt$ تعریف شده باشد. اگر شمار !!{element}s در~$A$ متناهی باشد،
آشکار است که نمی‌توانیم تابعی داشته باشیم که بر بی‌نهایت !!{element}sی
مجموعهٔ~$\PosInt$ تعریف شده باشد و همهٔ
!!{element}sی مجموعهٔ~$A$ را به‌عنوان مقدار بگیرد، اما هیچ‌گاه مقداری را
دوبار نگیرد. در این حالت، یعنی هنگامی که فهرست بی‌تکرار متناهی است، باید
دامنه‌ای متفاوت برای~$f$ برگزینیم که شمار !!{element}sی آن متناهی باشد.
نداشتن تکرار یعنی $f$ باید !!{injective} باشد. چون این تابع
!!{surjective} نیز هست، در پی !!a{bijection} میان یک مجموعهٔ متناهی
$\{1, \dots, n\}$ یا $\PosInt$ و~$A$ هستیم.

\begin{prop}\ollabel{prop:enum-bij}
اگر $f\colon \PosInt \to A$ !!{surjective} باشد (یعنی برشماریِ~$A$
باشد)، !!a{bijection} $g\colon Z \to A$ وجود دارد که در آن $Z$ یا
~$\PosInt$ است یا $\{1, \dots, n\}$ برای یک~$n \in \PosInt$.
\end{prop}

\begin{proof}
  تابع $g$ را به‌صورت بازگشتی تعریف می‌کنیم: قرار دهید $g(1) = f(1)$.
  اگر $g(i)$ پیش‌تر تعریف شده باشد، $g(i+1)$ را نخستین مقدار از $f(1)$،
  $f(2)$، \dots{} قرار دهید که پیش‌تر در میان $g(1)$، \dots، $g(i)$
  نباشد، البته اگر چنین مقداری وجود داشته باشد. اگر شمار !!{element}sی~$A$
  دقیقاً $n$ باشد، آنگاه $g(1)$، \dots، $g(n)$ همگی تعریف شده‌اند
  و بنابراین تابعی $g\colon \{1, \dots, n\}
  \to A$ تعریف کرده‌ایم. اگر شمار !!{element}sی~$A$ نامتناهی باشد، آنگاه برای هر $i$ باید
  !!a{element} از~$A$ در برشماریِ $f(1)$، $f(2)$، \dots وجود داشته باشد
  که پیش‌تر در میان $g(1)$، \dots، $g(i)$ نباشد. در این حالت، تابعی
  $g\colon \PosInt \to A$ تعریف کرده‌ایم.

  تابع $g$ !!{surjective} است، زیرا هر عضو از~$A$ در میان $f(1)$،
  $f(2)$، \dots{} قرار دارد (چون $f$ !!{surjective} است) و در نتیجه سرانجام
  مقدارِ~$g(i)$ برای یک~$i$ خواهد شد. این تابع همچنین !!{injective} است،
  زیرا اگر برای $j < i$ داشتیم $g(j) = g(i)$، آنگاه $g(i)$ پیش‌تر
  در میان $g(1)$، \dots، $g(i-1)$ قرار می‌داشت، برخلاف شیوه‌ای که~$g$
  را تعریف کردیم.
\end{proof}

\begin{cor}\ollabel{cor:enum-nat-bij}
مجموعهٔ $A$ !!{enumerable} است اگر و تنها اگر تهی باشد یا !!a{bijection}
$f\colon N \to A$ وجود داشته باشد که در آن یا $N = \Nat$ یا
$N = \{0, \dots, n\}$ برای یک $n \in \Nat$.
\end{cor}

\begin{proof}
$A$ !!{enumerable} است اگر و تنها اگر $A$ تهی باشد یا تابعی
!!a{surjective} به‌صورت $f\colon \PosInt \to A$ وجود داشته باشد. بنا بر
\olref{prop:enum-bij}، مورد اخیر برقرار است اگر و تنها اگر تابعِ
!!a{bijective}~$f\colon Z \to A$ وجود داشته باشد که در آن $Z =
\PosInt$ یا $Z = \{1, \dots, n\}$ برای یک $n \in \PosInt$. با همان
استدلالِ برهان \olref{prop:enum-shift}، این نیز به نوبهٔ خود برقرار است
اگر و تنها اگر !!a{bijection} $g\colon N \to A$ وجود داشته باشد که در آن
یا $N = \Nat$ یا $N = \{0, \dots, n-1\}$.
\end{proof}

\begin{prob}
  بنا بر \olref[sfr][siz][enm]{defn:enumerable}، مجموعهٔ $A$ شمارا است
  اگر و تنها اگر $A = \emptyset$ یا تابعی !!a{surjective} به‌صورت $f\colon
  \PosInt \to A$ وجود داشته باشد. همچنین می‌توان اصطلاح
  ``مجموعهٔ !!{enumerable}'' را به‌طور دقیق چنین تعریف کرد: یک مجموعه شمارا
  است اگر و تنها اگر تابعی !!a{injective} به‌صورت
  $g\colon A \to \PosInt$ وجود داشته باشد. نشان دهید که این دو تعریف
  هم‌ارزند؛ یعنی نشان دهید که تابعی !!a{injective} به‌صورت
  $g\colon A \to \PosInt$ وجود دارد اگر و تنها اگر یا $A = \emptyset$
  یا تابعی !!a{surjective}~$f\colon \PosInt \to A$ وجود داشته باشد.
\end{prob}

\end{document}
